\section{All four proved conductor-sixteen residues in the actual arithmetic window}
\label{sec:conductor16-four-residue-window}

\paragraph{Status and dependency boundary.}
This is a conditional-interface proof.  It uses the original strict
theorems AC-G142, AC-G144 and AC-G147, and the established-sector
comparison AC-G146.  Their exact proof and review status is recorded
separately.  The new class input is AC-G147-class, applied only through
AC-G147.  No remaining index bound or mixed realization is proved here.

Put $p=269$, $K=\mathbb Q(\omega)$, $t^3=p$, with $t$ real, and
let $\chi$ be the cubic character of $K(t)/K$.
Let $\phi$ be the primitive odd character modulo sixteen with
$\phi(13)=i$, $\phi(5)=-i$, and $\ker\phi=\{1,7\}$.
Write $\psi_4=\chi_{-4}|_{G_K}$, and use the actual cyclotomic
character $c_p$.  Both signs of $\phi$ and both nontrivial cubic
labels are always included, so changing the chosen roots of unity
does not change the finite label set.

Fix any source datum, any $0<\varepsilon<1/256$, its allowed page,
and its prescribed pair $q,d$.  Retain the actual fields
\[
 \begin{aligned}
 F_u&=\mathbb Q(\zeta_{2^{h(u)+2}},\zeta_7,
                  \sqrt{-q},\sqrt{-d},\omega),\\
 H_u&=F_u(t),\qquad G_u=\operatorname{Gal}(H_u/K),\\
 u^{2+8\varepsilon}&\leq2^{h(u)}\leq2u^{2+8\varepsilon},
 \qquad D_u=|G_u|,\qquad T=\lfloor u^{399}\rfloor.
 \end{aligned}
 \tag{CWF.1}
\]
All coefficient rings $O_\ell$, fields $E_\ell$, degrees $b_\ell$,
strict conditions and index normalizations are those of
(FC.1)--(FC.2), (EU.2) and (TC.1).  In particular the coefficient-prime
range remains unrestricted.

For $u\geq2$ define the following disjoint label sets:
\[
 \begin{aligned}
 Z_{16}^{+}&=\{(p,\chi^{-j}\phi^{-\sigma},e):
       j=1,2,\ \sigma=\pm1,\\
 &\hspace{2em}1\leq e<T,\ e\equiv-1\pmod{67}\},\\
 Z_4&=\{(p,\chi^{-j}\psi_4,e):
       j=1,2,\ 1\leq e<T,\ e\equiv-1\pmod{134}\},\\
 Z_{\rm cub}&=\{(p,\chi^j,e):
       j=0,1,2,\ 1\leq e<T,\ e\equiv0,267\pmod{268}\},\\
 Z_3&=\{(3,1,e):1\leq e<T\},\qquad
 Z^+=Z_{16}^{+}\sqcup Z_4\sqcup Z_{\rm cub}\sqcup Z_3.
 \end{aligned}
 \tag{CWF.2}
\]
Let $I_{399}^{\rm est+}$ be exactly (TC.1) with the labels $Z^+$
omitted.  Every other prime, finite character and positive Tate
weight remains.  Denote the old removed set of AC-G146 by $Z^{\rm est}$
and its remaining sum by $I_{399}^{\rm est}$.

\begin{theorem}[AC-G148-window: the four-residue actual-label transport]
\label{AC-G148-window}
On every prescribed host (CWF.1), all labels in (CWF.2) occur and
their original-coefficient strict groups (FC.2) vanish.
Their exact disjoint counts are
\[
 \begin{aligned}
 |Z_{16}^{+}|&=4\lfloor T/67\rfloor,&
 |Z_4|&=2\lfloor T/134\rfloor,\\
 |Z_{\rm cub}|&=3\bigl(\lfloor(T-1)/268\rfloor
                            +\lfloor T/268\rfloor\bigr),&
 |Z_3|&=T-1.
 \end{aligned}
 \tag{CWF.3}
\]
Relative to AC-G146, the additional omitted labels number exactly
\[
 |Z^+\setminus Z^{\rm est}|
     =4\left\lfloor\frac{T+67}{134}\right\rfloor.
 \tag{CWF.4}
\]
Write $C_u$ for the sum of the actual nonnegative normalized
positive-index costs over $Z^+$.  Then
\[
 \begin{aligned}
 I_{399}&=I_{399}^{\rm est+}+C_u,\qquad
 I_{399}^{\rm est+}\leq I_{399}^{\rm est}\leq I_{399},\\
 0\leq C_u&\leq
 8D_u^2T^2\log(8TD_u)+\tfrac12D_u\mathcal R_uT^2
                                      +2D_ur_uc_0T\\
 &=O_\varepsilon\bigl(u^{802+16\varepsilon}(\log u)^2\bigr).
 \end{aligned}
 \tag{CWF.5}
\]
The first equality also makes sense for nonnegative extended sums;
no subtraction of two infinite quantities is required.
There is only one full (FL.4) majorant for the entire set $Z^+$,
including the retained coefficient-three labels.

The AC-K35 common-window premise, and its AC-G146 version, may
therefore use $I_{399}^{\rm est+}$ with all the same simultaneous
quantifiers and positive relative slack.  The old premise implies
the new one with unchanged parameters.  The new premise implies
the full old one with the same $W,J$, datum, page and pair, after
only increasing the starting threshold and halving the slack.
The unchanged mixed theta/Selmer and analytic realization input
$\mathcal R_{\rm cub}$ is still required for the conditional
AC-K35 conclusion.  Its truth and the remaining numerical bound
are not asserted.
\end{theorem}

\begin{proof}
\proofstep{1. Actual finite characters and original coefficient rings.}
The page lower bound and $u\geq2$ imply $h(u)\geq2$.
Consequently $\mathbb Q(\zeta_{16})\subset F_u$ and
$K(t)\subset H_u$.  The character $\phi$ and the cubic character
both factor through $G_u$.  Since $F_u/\mathbb Q$ is abelian
and the only nontrivial abelian subfield of $K(t)$ is $K$,
$F_u\cap K(t)=K$.  The cubic subgroup distinguishes $j$ and
the quartic quotient distinguishes $\sigma$.
The quartic, quadratic and trivial two-power factors are distinct,
so the four sets in (CWF.2) are disjoint.  The coefficient-three
set is disjoint for its coefficient prime.  The other finite labels
and their occurrence are exactly those already verified in AC-G146.

The original ring $O_p$ contains the character values, since
the exponent $3\cdot2^{h(u)}$ of the actual group contains
three and four.  The original $b_p$ is retained.
For a conductor-sixteen label the coefficient is literally
\[
 A_{j,\sigma,e}=(O_p\otimes_{\mathbb Z_p}\mathbb Q_p/\mathbb Z_p)
              (\chi^{-j}\phi^{-\sigma})(-e)
       =(E_p/O_p)(\chi^{-j}\phi^{-\sigma}c_p^{-e}).
 \tag{CWF.6}
\]
Apply AC-G142 at residue $267$, AC-G144 at residue $133$,
and AC-G147 at residues $66$ and $200$, always with the additional
finite exponent zero and the continuous character one.
Each theorem holds on this given $O_p$, including when ramified.
The order-four cyclotomic character in the new residual Kummer
field is not inserted into (CWF.6) or the finite host-label list.
For $Z_4,Z_{\rm cub},Z_3$, retain the exact original-coefficient
strict assertions and transports proved in AC-G146.

\proofstep{2. Full local conditions and exact counts.}
The three conductor-sixteen strict theorems impose zero in full
$H^1(G_{K_v},A)$ for $v\mid p$ and zero in $H^1(I_v,A)$ at
every other finite place, exactly as in (FC.2).
The finite-set presentation agrees with the actual $\Sigma_p$:
at an added away-$p$ place outside $2,3,p$, (CWF.6) is unramified.
A strict inertia restriction is a coboundary, hence is zero
because inertia acts trivially.  The cocycle identity then gives
zero on the closed normal subgroup generated by all added inertia.
Both the action and the cocycle descend to the smaller group.
Inflation is injective and preserves the remaining local conditions.
Thus the strict kernels agree, without dropping a full local place.
The same exact comparison for the other sets is retained from AC-G146,
including its coefficient-three argument without averaging by $3$.

The four conductor-sixteen progressions are exactly
$e=67n-1$, $n\geq1$.  The inequalities $1\leq67n-1<T$
are equivalent to $1\leq n\leq\lfloor T/67\rfloor$.
The old two progressions have even $n$; the new two have odd $n$,
or $e=134m+66$, $m\geq0$.  Hence
\[
 \lfloor T/134\rfloor+
 \left\lfloor\frac{T+67}{134}\right\rfloor
                         =\lfloor T/67\rfloor.
 \tag{CWF.7}
\]
Multiplication by four proves the first count and (CWF.4).
The two conductor-four labels retain $e=134n-1$; the pure-cubic
and coefficient-three counts are unchanged from (EW.2).
This proves all of (CWF.3).  Exactly these strict terms in (FC.15)
are zero; no assertion about an individual index follows from that alone.

\proofstep{3. One error for the complete union.}
For a removed label $(\ell,\eta,e)$ use the original dual complex
$C=R\Gamma(G_{K,\Sigma_\ell},O_\ell(\eta^{-1})(e+1))$,
with the codifferent convention of (FL.2).
The retained determinant inequality (EU.2) and strict zero give
\[
 \frac{\log\#k_{O_\ell}}{b_\ell}\max(0,i_{\ell,\eta,e})
 \leq\frac{\log\#H^2(C)}{b_\ell}
 =\frac{|\log\#H^2(C)-\log\#\mathcal S_{\ell,\eta}(e)|}{b_\ell}.
 \tag{CWF.8}
\]
This equality uses only $\mathcal S=0$; it does not identify
strict cohomology with $H^2$.
For a fixed weight, sum (CWF.8) over all removed labels at both
coefficient primes.  It is one subsum of the nonnegative absolute
comparison (FL.4), so its bound is
\[
        16eD_u^2\log(8eD_u)+D_ue\mathcal R_u+2D_ur_uc_0.
 \tag{CWF.9}
\]
Sum this once over $1\leq e<T$, using
$\sum e\leq T^2/2$, $\log(8eD_u)\leq\log(8TD_u)$ and
$T-1\leq T$.  This proves (CWF.5) without adding a second full
error for the new residues or for any old family.
Nonnegativity gives additivity and both monotonicity inequalities.

The unchanged actual-host estimates are
$D_u\leq288u^{2+8\varepsilon}$, $r_u\leq12$ and
$\mathcal R_u=O_\varepsilon(u^{\varepsilon/5}(\log u)^2)$,
uniformly over the prescribed page and pair.
With $T\leq u^{399}$ the first error exponent is
$802+16\varepsilon$, the second is $800+41\varepsilon/5$,
and the third is smaller.  Thus the entire bound is
$o(\exp(u^\varepsilon/2))$ for every fixed positive $\varepsilon$.

\proofstep{4. Preserve the single common window.}
Fix the same datum, page and $0<\varepsilon<1/256$.
Suppose the new premise supplies $J\geq1$, $X_0$ and $s\in(0,1)$
such that for every $X\geq X_0$ and every prescribed
$\operatorname{WinPair}(\varepsilon,X^{1/J},q,d)$ it supplies one
$W$ with $X^{1/J}\leq W$, $W^2\leq X$, and
\[
 I_{399}^{\rm est+}(H_u,u)\leq(1-s)\exp(u^\varepsilon/2)
                           \quad(W\leq u\leq W^2).
 \tag{CWF.10}
\]
Choose $U_0\geq2$ so that the uniform error in (CWF.5) is at most
$(s/2)\exp(u^\varepsilon/2)$ for all $u\geq U_0$.
It depends only on the same datum, page, $\varepsilon$ and slack,
not on a new pair, coefficient, character, Tate weight or window.
Increase $X_0$ to $\max(X_0,U_0^J)$.  The very same $W$ then gives
\[
 I_{399}(H_u,u)\leq(1-s/2)\exp(u^\varepsilon/2)
                              \quad(W\leq u\leq W^2).
 \tag{CWF.11}
\]
For any fixed finite simultaneous source family take the maximum
of its uniform error constants, retaining this one $W$.
No intersection of independently chosen windows is used.
Conversely either old premise implies (CWF.10) by nonnegativity
with all its original parameters; the little-$o$ versions likewise
agree.  The actual $R.t$, source host, family construction and
every retained $\mathcal R_{\rm cub}$ clause are unchanged.
This is exactly the conditional implication claimed.
\end{proof}

\paragraph{Existing-gate and geometric scope.}
The AC-G146 unknown sum still includes the four conductor-sixteen
labels at $e=66,200\bmod268$.  For $T\geq67$, (CWF.4) proves
proper inclusion of the omitted label sets, on the same actual hosts.
It does not assert a strict numerical inequality of sums.
The conductor-four, pure-cubic and coefficient-three omissions
are retained, not counted as new statements.
The unlisted all-prime/character/weight bound, mixed theta/Selmer
and analytic realization, and H04.4 remain open.
The coefficients remain the original rank-one Artin--Tate systems
on the same arithmetic curve and finite cover.  SGA~1,
Expos\'e~V, \S\S4 and 7 supplies only the previously retained
finite-cover dictionary.  No new motive, reciprocity law,
primary-source proof or Lean execution is asserted.
